\section{Introduction}
\label{sec:intro}
Spatial intelligence, the ability to understand geometric relations among objects and their 3D environments, is essential for vision-language models (VLMs) to operate in the physical world and represents a key step toward artificial general intelligence (AGI), where models are expected to perceive, reason, and make decisions in 3D space as humans do. 
% 
Such capability is crucial for real-world applications, including embodied robotics \citep{driess2023palme,brohan2023rt2}, AR/VR perception \citep{newcombe2011kinectfusion}, and autonomous driving \citep{geiger2012kitti,caesar2020nuscenes}. 
% 
However, unlike human perception, which naturally integrates visual cues into coherent 3D understanding, current VLMs are primarily trained on passive 2D visual-text corpora, with limited explicit 3D supervision or embodied experience \citep{radford2021clip,alayrac2022flamingo,liu2023llava}.
% 
This creates a fundamental semantic-to-geometric gap: while VLMs excel at probabilistic and qualitative semantic inference, their reasoning is often mediated by lossy semantic representations that fail to faithfully capture high-fidelity geometry, leaving them susceptible to textual patterns and semantic priors rather than grounded 3D geometric evidence \citep{gca,testset_stress}.


Recent advances in agentic VLMs substantially push the boundary of spatial understanding by augmenting VLMs with external tools, executable programs, and explicit geometric structure. 
% 
For example, VADAR \citep{vadar} dynamically constructs a Python API and synthesizes programs for 3D spatial reasoning; SpaceTools \citep{spacetools} trains VLMs to coordinate multiple vision and robotic tools through interactive reinforcement learning. %. %, and GCA \citep{gca} constrains the reasoning process with formal reference-frame and objective constraints before performing deterministic geometric computation.
% 
However, despite their strong performance, these methods still largely focus on static images or isolated visual observations, which remains far from the goal of real-world spatial intelligence: \textit{the real 3D world is hidden, evolving, and continuously projected into streams of 2D observations}. 
% 
Reasoning from isolated 2D views alone makes it fundamentally challenging to maintain persistent object states, integrate evidence across viewpoints and time, and build a coherent understanding of the underlying 3D scene.

To move beyond static and stateless spatial reasoning, we introduce \textbf{\textsc{S-Agent}}, a \textbf{S}patial tool-use agentic paradigm for understanding and reasoning over continuous multi-view images and videos. 
% 
Our key motivation is that the missing ingredient for video-based spatial intelligence is not merely stronger 2D/3D visual recognition, but a reasoning mechanism that can accumulate spatial evidence along both spatial and temporal dimensions. 
% 
Specifically, in continuous multi-view and video settings, each frame is only a partial and transient observation of the scene, while the key to spatial intelligence is to connect these observations into a spatially structured and temporally persistent understanding of the underlying 3D world.
% 
Rather than asking VLMs to implicitly internalize this entire process, our \textsc{S-Agent} casts the VLM as a semantic planner that decides what evidence is needed, while spatial tools/experts and temporal memory provide continuous and explicit 3D awareness of the specific scene, ranging from low-level 2D/3D evidence (e.g., object grounding, depth information) to high-level spatial knowledge (e.g., orientations, relationships).
% 
This separation enables the agent to reason from accumulated evidence instead of isolated visual impressions, extending existing spatial agent methods toward stateful, temporally grounded understanding of evolving scenes.
% 
% Unlike prior spatial agents that solve each query from isolated observations, \textsc{S-Agent} maintains an evolving scene state and reasons over accumulated evidence across both viewpoints and time.

Motivated by this perspective, \textsc{S-Agent} is designed as a VLM-orchestrated spatio-temporal reasoning framework: it progressively aggregates spatial evidence from fragmented 2D observations into structured 3D scene knowledge, while persistently accumulating temporal evidence across frames and reasoning iterations. 
% 
(1) For the \textit{spatial dimension}, \textsc{S-Agent} follows a hierarchical understanding process. 
At the first level, \textbf{2D perception tools} ground objects and regions in individual frames, establishing object-centric visual facts for subsequent reasoning. 
At the second level, \textbf{multi-view 3D tools} enrich these grounded entities with geometric cues (e.g., depth, 3D coordinates, and camera poses), allowing evidence from different viewpoints to be integrated beyond the original image plane. 
At the third level, \textbf{specialized spatial experts} aggregate these geometric signals into higher-level spatial knowledge (e.g., object counts, physical measurements, orientations, and relative positions). 
% 
(2) For the \textit{temporal dimension}, \textsc{S-Agent} maintains memory over the evolving reasoning process: \textbf{Scene Memory} tracks grounded entities across frames to preserve object identity and suppress duplicate evidence, while \textbf{Agent Memory} stores accumulated tool observations and intermediate reasoning traces for iterative refinement. 
In this way, \textsc{S-Agent} turns video spatial reasoning from disconnected frame-level prediction into evidence accumulation over an evolving 3D scene.


Comprehensive experiments on multi-image benchmarks (MMSI-Bench~\citep{mmsi_bench} and ViewSpatial-Bench~\citep{viewspatial_bench}) and video spatial reasoning benchmarks (ReVSI~\citep{revsi} and VSI-SUPER~\citep{cambrian_s}) validate the robustness and generalizability of our approach.
% 
(1) \textbf{Zero-shot setting}. We directly instantiate \textsc{S-Agent} with both open-source models (e.g., Qwen3) and closed-source APIs (e.g., Gemini and GPT) in a training-free manner. 
Simply and directly applying the \textsc{S-Agent} framework consistently improves the spatial reasoning ability of these VLMs, improving over GPT-5.4 by 4.5\% on MMSI-Bench. 
% 
(2) \textbf{Training setting}. Beyond inference-time improvement, we further construct a spatial-instruction dataset \textsc{S-300K} from zero-shot \textsc{S-Agent} trajectories on the SenseNova-SI-800K \citep{sensenova_si} training set (which is fully disjoint from all evaluation benchmarks) and use it to perform supervised fine-tuning on Qwen3-VL-8B, resulting in \textsc{S-Agent-8B}. 
% 
Compared with direct Qwen3-VL-8B inference, \textsc{S-Agent-8B} achieves a 10.5\% improvement on MMSI-Bench, improving accuracy from 31.1\% to 41.6\%, and performs comparably to advanced closed-source models such as GPT-5.4 and Gemini 3 Pro across multiple benchmarks. 
% 
These results show that \textsc{S-Agent} is not only an effective training-free inference framework, but also a scalable paradigm for building compact spatially capable agents.
